\section{One reference level controls every weighted slow-band tail}
\label{sec:weighted-slow-band}

The actual nested PW ranges give a shorter uniform tail
argument than a growing-metric Gram estimate.  It works
in the single physical Hilbert space, with its specified
weight and orthogonal projections.  It therefore retains
the boundary jets that occur at positive polynomial weight.

\paragraph{MA-SB-005 (the actual slow-band limit at every polynomial weight).}
Fix the rational trivial sector,
$0<\alpha<1/2$, $s\ge0$, and
\[
 \begin{aligned}
 H&=L^2(\mathbb R,w(x)\,dx),&
 w(x)&=e^{2\alpha|x|}(1+x^2)^{s/2},\\
 P_c&=M_{\mathbf1_{[c,\infty)}},& H_c&=P_cH,\\
 J_c&=J_{e^{-c}},& Z_c&=H_c\ominus J_c.
 \end{aligned}
 \tag{MA-SB.50}
\]
Here $J_Y$ is the actual source PW theta-image closure,
as in MA-GP-002/010 and MA-GW.  Write $P_c^Z$ for
its actual complement projection, extended by zero
on $H_c^\perp$, so that
$Q_Y^{\rm glob}=P_c-P_c^Z$ when $c=-\log Y$.
Let $g\in C_c^\infty((0,\infty))$ and choose
$0<a<b$ with $\operatorname{supp}g\subset[a,b]$.
The single fixed reference operator is
\[
 K_{0,g}=C_gP_0^Z\in\mathcal S_1(H),\qquad
 \Omega_g(R)=\|P_RK_{0,g}\|_1
                       \longrightarrow0\quad(R\to\infty).
 \tag{MA-SB.51}
\]
Its trace class is MA-GW-003 at source level $Y=1$.
For every $c\le0$ and $R>b$ one has the exact identity
and uniform bound
\[
 \begin{aligned}
 P_RC_gP_c^Z&=P_RC_gP_0^ZP_0P_c^Z,\\
 \|P_RC_gP_c^Z\|_1&\le\Omega_g(R).
 \end{aligned}
 \tag{MA-SB.52}
\]
The same estimate holds with any multiplication
profile of norm at most one supported on $[R,\infty)$
in place of the leftmost projection.

Put $d=1/2-\alpha$, $u=1/2+\alpha$, and let
\[
 \begin{aligned}
 b_s&=\#\{j\in\mathbb Z_{\ge0}:2j<s-1\},\\
 r_\rho^+&=
 \begin{cases}
 m_\rho,&\operatorname{Re}\rho<u,\\
 \min(m_\rho,b_s),&\operatorname{Re}\rho=u,\\
 0,&\operatorname{Re}\rho>u,
 \end{cases}\\
 \tau_{+,s}(g)&=\sum_\rho r_\rho^+G_0(\rho-1/2),\qquad
 G_0(z)=\int g(x)e^{zx}\,dx.
 \end{aligned}
 \tag{MA-SB.53}
\]
Take a source cutoff $0\le\rho_{X,\delta}\le1$
supported on $[1/X,X]$ and equal to one on
$[1/X+\delta,X-\delta]$.
For
$Y\ge1$, $X\ge2Y$, $0<\delta\le1/(4X)$, and
$\log(X-\delta)>b$, the actual trace satisfies
\[
 \begin{split}
 &\left|\operatorname{Tr}_H
       (M_{\rho_{X,\delta}}C_gQ_Y^{\rm glob})
                     +\tau_{+,s}(g)\right|\\
 &\hspace{20mm}\le\Omega_g(\log(X-\delta))
                         \le\Omega_g(\log(X/2)).
 \end{split}
 \tag{MA-SB.54}
\]
Thus every source-thin path with $X,Y\to\infty$ and
$Y/X\to0$ has actual limit $-\tau_{+,s}(g)$.
The estimate is uniform for $1\le Y\le X/2$.
It covers every $s\ge0$, in particular the draft's
$s\ge3$, with precisely the upper-edge lengths
in (MA-SB.53).

\begin{proof}
\emph{An exact orthogonality identity in the actual norm.}
The smooth source range at $Y=1$ is contained in
the smooth source range at every $Y\ge1$.
Equivalently MA-GP-002 gives the supported-ideal
inclusion at the two endpoints.  Taking closures
in this same $H$ gives
\[
                         J_0\subset J_c\qquad(c\le0).
 \tag{MA-SB.55}
\]
No equivalence of norms depending on $c$ is used.
If $z\in Z_c$ and $j\in J_0$, then $P_0j=j$ and
\[
              \langle P_0z,j\rangle_H
                 =\langle z,j\rangle_H=0.
\]
As $P_0z\in H_0$, it follows that $P_0z\in Z_0$.
Consequently the actual orthogonal projections obey
\[
                  P_0P_c^Z=P_0^ZP_0P_c^Z,
                     \qquad\|P_0P_c^Z\|\le1.
 \tag{MA-SB.56}
\]
The final norm bound is in the same physical weighted
space on both sides.

\emph{Causality puts every far right tail through the reference level.}
Convolution $C_g$ is bounded on $H$ by the translation
estimate of MA-WP.6--.7.  If an input is supported
on $(-\infty,0)$, its output under $C_g$ is supported
on $(-\infty,b]$.  This follows directly from
$(C_gf)(x)=\int g(t)f(x-t)dt$, first on compact
smooth inputs and then by boundedness and density.
Hence $P_RC_g=P_RC_gP_0$ for $R>b$.
Insert (MA-SB.56) to obtain
\[
 \begin{split}
 P_RC_gP_c^Z
   &=P_RC_gP_0P_c^Z\\
   &=P_RC_gP_0^ZP_0P_c^Z .
 \end{split}
 \tag{MA-SB.57}
\]
The reference factor $K_{0,g}$ is trace class by
MA-GW-003.  The trace-ideal inequality and the
contraction in (MA-SB.56) give (MA-SB.52).
The profile variant factors a norm-one multiplication
operator through $P_R$.
Strong convergence $P_R\to0$ as $R\to\infty$,
followed by finite-rank approximation of $K_{0,g}$,
proves (MA-SB.51).  Monotonicity of $\Omega_g$ follows
from $P_{R'}P_R=P_{R'}$ when $R'\ge R$.

This tail factorization in fact needs only the
reference-level ambient lift.  Each full
$C_gP_c^Z$ is then trace class: its right tail
is (MA-SB.57), and its output on $[c,R]$ is
a compact-position smooth convolution kernel,
with bounded compact weight factors and support
projections.  Such a kernel is trace class by
the compact Fourier-series argument already
used in MA-GP-008.  Causality gives zero output
below $c$.  This also proves the statement
without any uniform estimate of the full
trace norm of $C_gP_c^Z$.

\emph{The actual trace and the source cutoff.}
MA-GW-002--003 identify the actual quotient and
ambient trace as $\tau_{+,s}(g)$, independent of
$c$.  The retained boundary lengths are those
in (MA-SB.53); they have not been replaced by
the ordinary $s=0$ Hardy divisor.
For every fixed compact position profile $\rho$,
the actual decomposition
$Q_Y^{\rm glob}=P_c-P_c^Z$ gives
\[
 \begin{aligned}
 \operatorname{Tr}(M_\rho C_gQ_Y^{\rm glob})
      &=-\operatorname{Tr}(M_\rho C_gP_c^Z),\\
 \operatorname{Tr}(M_\rho C_gQ_Y^{\rm glob})
                       +\tau_{+,s}(g)
      &=\operatorname{Tr}((1-M_\rho)C_gP_c^Z).
 \end{aligned}
 \tag{MA-SB.58}
\]
All products are trace class before taking traces.
The first equality uses the compact-position
kernel trace $\operatorname{Tr}(M_\rho C_gP_c)=0$
proved in MA-GW.44; its diagonal is zero since
$g(0)=0$.

Now take the source profile and $c=-\log Y$.
Convolution is causal, so $C_gP_c^Z$ has range
in $H_c$.  On that half-line,
$e^x\ge1/Y\ge2/X>1/X+\delta$.
The entire lower transition is therefore absent.
The profile is one through
$R=\log(X-\delta)$.
Consequently the multiplier
$(1-\rho(e^x))\mathbf1_{[c,\infty)}(x)$ has norm
at most one and is supported on $[R,\infty)$.
Apply the profile version of (MA-SB.52) to the
second line of (MA-SB.58).  This proves the first
bound of (MA-SB.54).  The second follows from
$X-\delta\ge X/2$ and monotonicity of $\Omega_g$.
The hypotheses are eventually satisfied on every
stated slow-band path with the source's thin
boundary condition $\delta X\to0$.
\end{proof}

\paragraph{Actual local comparison and both boundary deficiencies.}
Use the common annular profile
$a_X(t)=A(t/X)-A(Xt)$ from MA-LA-014,
with $A=1$ on $[0,1]$, zero on $[2,\infty)$,
and values in $[0,1]$.
For $X\ge2Y$ its lower term vanishes on $H_c$,
and its upper term is one through $x=\log X$.
The same proof gives error at most
$\Omega_g(\log X)$ and the same global limit.
By MA-GP-010 this is also the actual finite-$S$
invariant global projection, with $S'=S$.

Choose $S$ to contain every prime-power contribution
to the fixed $h(t)=g(\log t)$.
The actual local quotient operator
$T^S_{X,Y,X^{-1}}(h)$ of MA-LA-014, with the
same annular profile, has limit $A_0(h)+B_S(h)$.
This local operator and its Hilbert norm are the
ones already constructed there; no polynomial
global weight is inserted into its definition.
The unconditional smooth Meyer trace identity
gives $A_0+B_S-P(g)=-\tau_{\rm all}(g)$, where
$P(g)=G_0(1/2)+G_0(-1/2)$.
The two actual limits thus give
\[
 \begin{split}
 &\lim_{\substack{X,Y\to\infty\\Y/X\to0}}
 \left\{\operatorname{Tr}(M_{a_X}C_gQ_Y^{\rm glob})
       -\bigl(\operatorname{Tr}T^S_{X,Y,X^{-1}}(h)-P(g)\bigr)\right\}\\
 &\quad=
   \sum_{\operatorname{Re}\rho>u}m_\rho G_0(\rho-1/2)\\
 &\qquad+\sum_{\operatorname{Re}\rho=u}
      \bigl(m_\rho-\min(m_\rho,b_s)\bigr)G_0(\rho-1/2).
 \end{split}
 \tag{MA-SB.59}
\]
Thus at the upper edge only the unretained
length remains, rather than the full multiplicity
that occurs at $s=0$.
Relative to the actual completed full-strip
trace of MA-WP, the difference is
\[
 \begin{split}
 &\lim_{\substack{X,Y\to\infty\\Y/X\to0}}
       \operatorname{Tr}(M_{a_X}C_gQ_Y^{\rm glob})
                       -\bigl(-\tau_{{\rm strip},s}(g)\bigr)\\
 &\quad=-\sum_{\operatorname{Re}\rho<d}
                         m_\rho G_0(\rho-1/2)\\
 &\qquad-\sum_{\operatorname{Re}\rho=d}
      \bigl(m_\rho-\min(m_\rho,b_s)\bigr)G_0(\rho-1/2).
 \end{split}
 \tag{MA-SB.60}
\]
These series are absolutely convergent.
They express the differences of already constructed
operator traces.  No omitted zero or omitted edge
jet is asserted to exist.
The positive-length divisor of each nonzero deficiency
also falls under the exact separation lemma MA-JP-004.

\paragraph{Abstract character scope and the retained Gram factors.}
The proof of (MA-SB.52) is valid for any locally
bounded positive weight with locally bounded inverse,
bounded $C_g$, actual closed supported invariant subspaces
$J_c$ satisfying $J_0\subset J_c$ for $c\le0$,
and the single actual hypothesis
$C_gP_0^Z\in\mathcal S_1$.
If their actual quotient trace is independent of
$c$, the same source-cutoff limit follows.
This gives an exact criterion for the character
families of MA-CGP: their supported ranges and
nesting are proved, and a positive-weight reference
lift with its boundary multiplicities is the only
additional hypothesis needed to apply this argument.
For $s=0$ that hypothesis is already MA-CGP-004,
so MA-SB-004 is recovered directly.
An application at $s>0$ uses a proved weighted
character lift; it does not treat a Blaschke
multiplier as unitary in the stronger norm.

For every fixed $X,Y$ the global trace in this
theorem is still the actual Gram-resolvent trace
of MA-GP-008 (or MA-CGP.18 in its proved character
scope).  The reference-level factorization estimates
that actual trace; it does not discard either
normal-equation factor or replace the projection
by an additive Fourier cutoff.
Its regularizer may first be removed at fixed
data, or selected along the path to give a trace
error tending to zero, exactly as in MA-SB-003.
No uniform regularizer rate is asserted.

Only fixed-operator trace-ideal continuity and
orthogonality in the actual Hilbert metric enter
the uniform tail estimate.  The all-weight
finite-length modules and their boundary truncations
enter through MA-GW's separately proved actual
trace.  This retains the cote~1 trace-framework
placement of BR-TR-01/02 without adding an
arithmetic positivity or zero-location premise.
